\documentclass{article}
\usepackage{amsmath,amssymb,amsthm}
\newtheorem{theorem}{Theorem}
\title{Post-tightening: a directed full-grade comparison, not an isomorphism}
\date{16 November 2026}
\begin{document}\maketitle
\paragraph{Scope.}
This complementary candidate checks the full weighted constant-polynomial
grade interface for guarded incremental arena feedback. Operational syntax
is supplied by the arena/feedback manuscripts, not redefined here. Let
$P:A+U\rightsquigarrow B+U$ be guarded with full grade $C$ and let
$Q:B\rightsquigarrow Z$ have full grade $D$. Let $I$ be the unary unit
grade. Write $L_Q:B+U\rightsquigarrow Z+U$ for the typed case which executes
$Q$ at $B$ and is the zero-event identity at $U$, retaining the full case
grade $I+D$. No grading is restricted to realized bodies.

\section{The full grade map}
\begin{theorem}
There is a canonical arity-preserving weighted collection map
\[
 \phi_{C,D,U}:(R_UC)\circ D\longrightarrow
 R_U\bigl(C\circ(I+D)\bigr). \tag{1}
\]
It is injective on shapes, has bijective ordered external-position maps,
and its actual coefficient map at each shape is an isomorphism. It is
natural in all arity-preserving weighted collection maps $C,D$ and loop-label
maps $U\to U'$. These assertions do not say its full shape map is surjective.
\end{theorem}
\begin{proof}
A source shape is a nonempty finite $C$ invocation tree $t$, with internal
edges labelled by $U$, and one $D$ shape at each external leaf of $t$.
At each $C$ node, decorate an internal-edge slot with the $I$ alternative;
decorate an external-leaf slot with its specified $D$ shape. The new
invocation node has shape in $C\circ(I+D)$. Its $I$ output continues along
the original labelled internal edge, and every output of each $D$ decoration
is external. Retain the original invocation-node tree and all loop labels.
This gives a target shape on every source shape, whether used or unused.

External positions are the ordered pairs consisting of an original external
leaf and a position in its $D$ shape; both sides have exactly these positions.
In particular the backward position map of the corresponding polynomial
transformation is this bijection, not an arbitrary map collapsing leaves.
The source coefficient is the tensor of all $C$ node factors followed by the
$D$ leaf factors in its declared parenthesization. The target has the same
factors grouped nodewise, with additional unit factors at internal slots.
The canonical labelled-factor permutation, associators and unitors give an
actual invertible coefficient arrow. Symmetric monoidal coherence makes the
construction independent of presentation choices and natural in the actual
coefficient arrows.

The image consists exactly of target trees satisfying: every $I$-decorated
slot has its unique output attached to an internal $U$ edge; every output
of every $D$-decorated slot is external. Delete each internal $I$ decoration
and move the $D$ decorations to the external leaves of the remaining $C$
tree. This is the unique inverse on the image, proving injectivity. Nullary
$D$ decorations are retained at their slots, so this inversion does not lose
them merely because they contribute no output leaf. The same rule applies
to nullary $C$ nodes. Native shape maps preserve the alternative tags, ordered
positions and loop labels under their given image; hence both naturality
routes agree on complete shapes and on every labelled coefficient factor.
\end{proof}

\paragraph{Typed refinement, not reachable-only quotient.}
The image condition is an independently decidable sort constraint on every
formal target shape. It defines the maximal subcollection on which the
specified inverse reconstruction works, not a deletion by probability or
reachability. It may itself contain arbitrarily many unused shapes. Passing
to it would be an explicit typed refinement of the target grade; the native
comparison (1) does not silently replace the original full target by it.

\section{Actual operational comparison}
The operational configurations of $\operatorname{Fb}_U(P);Q$ consist of a
finite compatible tree of $P$ configurations and one $Q$ configuration at
each currently emitted $B$ port of that tree. Those of
$\operatorname{Fb}_U(P;L_Q)$ consist of a tree of $P;L_Q$ configurations:
each contains the same $P$ configuration, a $Q$ copy at each emitted $B$
port, and a trivial identity copy at each emitted $U$ port. The latter port
initializes the corresponding child invocation immediately. Removing these
trivial identity copies and flattening the stable path tags gives inverse
configuration functions, including initially unreachable configurations.
Use the repaired feedback compiler's label-qualified invocation edges:
$(p,u)$ on the left corresponds on the right to
$((p,\text{fixed }U\text{-identity port tag}),u)$.
Flattening and restoring these tags are global identifier functions, not
configuration-dependent renamings. Every $Q$-generated port has label in
$Z$, never in $U$, including at unreachable typed configurations. Thus the
same path correspondence applies to the entire occurring namespace.
Case identifiers are tagged by the identity versus $Q$ branch, and the
pure identity uses its fixed singleton port identifier independently of
input values. Stable path flattening retains these tags until the specified
identity-copy removal, giving consistent global identifier maps.
Initial $P$ configurations have no $U$ port; initial $Q$ copies cannot emit
$U$, so the right-hand feedback remains guarded.

Enabled primitive events on both sides are the same $P$ events and the same
$Q$-copy events, with identical normalized rows and actual context
coefficients. Their updates create the same child or continuation copies,
keep the same old ones, and commute with the flattening. Final output ports
have the same ordered stable paths. Thus this is an operational isomorphism
of the configuration/event/port data. At terminal configurations the full
body comparison is precisely (1): original recursive slots use $I$ and
original external slots use the actual terminal $Q$ shape.

Consequently, assuming the supplied arena composition and feedback
constructions, these data form an actual native arrow
\[
 \operatorname{Fb}_U(P);Q\longrightarrow
 \operatorname{Fb}_U(P;L_Q). \tag{2}
\]
Its configuration part is invertible and its primitive row maps are the
structural identity isomorphisms, but its full grade map is (1).
The inverse configuration function alone is not a native inverse. Every
complete selected policy transports through the operational event/scope
bijections, including zero and off-policy tables, without averaging or
rescheduling. This does not invert the full native grade comparison.

\section{A decisive reverse obstruction}
Take $U$ singleton, $C$ a singleton unary grade, and $D$ a singleton binary
grade, all with unit coefficients. Every finite $R_UC$ tree is a unary chain
with one external leaf. Hence every shape of $(R_UC)\circ D$ has arity two.
In the right-hand grade, a root $C$ node decorated by $D$ has two outputs.
Make one output external and attach the other by a $U$ edge to another
$C$ node decorated by $D$, whose two outputs are external. This is a valid
full formal target shape of arity three. There is no arity-preserving map
from the full right-hand collection to the left-hand collection. The shape
is unused by actual typed $P;L_Q$ executions, but remains part of their full
global grade. Thus the failure is not repaired by an operational state
isomorphism or by ignoring zero-weight witnesses.

\section{Coherence with further postcomposition}
Successive tightening must also retain unused case shapes. There is a
canonical directed map
\[
 \delta_{D,E}:I+(D\circ E)\longrightarrow
 (I+D)\circ(I+E): \tag{3}
\]
a $D\circ E$ shape chooses the $D$ alternative and $E$ at every position;
the unit alternative chooses $I$ then $I$. Its coefficients use the canonical
unit/tensor constraints. It is also the full grade map of the actual native case comparison
$L_{Q;R}\to L_Q;L_R$ when $Q$ has grade $D$ and $R$ grade $E$.
Its configuration map flattens the same continuation copies; on the identity
branch it inserts the two trivial identity copies. Tagged identifiers give
actual maps, and the coefficient arrows are those of (3). Applying guarded
feedback supplies the comparison used in the common-target diagram.
This is generally not a full-grade isomorphism:
the target also contains formal mixtures of $I$ and $E$ below a $D$ shape.

Starting with $((R_UC)\circ D)\circ E$, there are two routes to the common
full grade
\[
 R_U\bigl((C\circ(I+D))\circ(I+E)\bigr).
\]
One first applies $\phi_{C,D,U}\circ E$ and then
$\phi_{C\circ(I+D),E,U}$. The other uses source substitution associativity,
then $\phi_{C,D\circ E,U}$, then $R_U(C\circ\delta_{D,E})$ and target
associativity. They agree. At every original internal $U$ slot both insert
$I$ then $I$; at each original external slot both insert the specified $D$
shape and then the specified $E$ shapes at all its positions. They have the
same backward position bijection and the same actual map on every labelled
$C,D,E$ coefficient factor. Symmetric monoidal coherence proves equality
of those arrows, not merely equality of carrier functions.

For three or more postcompositions the same all-internal/all-external choice
rule proves the corresponding common-target diagrams: no slot changes its
sort along a route, and all coefficient maps are the same labelled-factor
coherence map. Operationally both routes flatten the identical family of
continuation configurations and primitive events, so the native diagrams
commute as well. This is directed full-grade coherence, not an identification
of all the intermediate target grades. For unit postcomposition, the fold $[1_I,1_I]:I+I\to I$,
followed by the substitution unitor, induces an actual directed grade map
\[
 R_U(C\circ(I+I))\longrightarrow R_UC.
\]
Composing it with $\phi_{C,I,U}$ is the original grade unitor
$(R_UC)\circ I\to R_UC$: both alternative tags are sent to the same
unary unit shape and all original $C$ nodes and loop labels remain intact.
At the arena level the typed case of two pure identities maps natively to
the pure identity, with per-configuration singleton-port bijections. Its two
branch-tagged global port identifiers may both map to the fixed identity
identifier, so this map is not asserted to be an operational isomorphism.
It yields a directed native unit comparison whose composite with tightening
is the original native unitor. On all formal shapes the analogous unit
diagrams with (3) agree by folding the same alternative tags; the actual
coefficient maps agree by monoidal unit coherence. No inverse on the full
right-hand shape collection or global identifier namespace is claimed.
\end{document}
